% =====================================================================
%  CHƯƠNG 1. BẢN CHẤT VÀ BA TRỤ CỘT NỀN TẢNG CỦA AI AGENT
% =====================================================================
\chapter{Bản chất và ba trụ cột nền tảng của AI Agent}

\section{Định nghĩa học thuật theo Google Cloud và McKinsey}
Google Cloud định nghĩa AI Agent là hệ thống phần mềm dùng trí tuệ nhân tạo để theo đuổi mục tiêu và hoàn thành nhiệm vụ
thay mặt người dùng; tác tử thể hiện khả năng suy luận, lập kế hoạch và ghi nhớ, đồng thời có một mức độ tự chủ để tự ra
quyết định, tự học và tự thích nghi \cite{google_agents}. McKinsey nhấn mạnh tác tử là hệ thống dựa trên mô hình nền
tảng, có thể \emph{thực thi} các quy trình nhiều bước phức tạp trong thế giới số: nhận một yêu cầu bằng ngôn ngữ tự
nhiên, tự chia nhỏ thành các bước, gọi công cụ và hoàn thành công việc thay vì chỉ trả lời \cite{mckinsey_agents}.

\begin{dinhnghia}{Định nghĩa làm việc dùng trong báo cáo}
\textbf{AI Agent} là hệ thống phần mềm vận hành bởi mô hình AI tạo sinh, có mức độ tự trị để chủ động theo đuổi một mục
tiêu: tự lập kế hoạch, tự duy trì ngữ cảnh và tự sử dụng công cụ bên ngoài để hoàn thành chuỗi tác vụ phức hợp thay mặt
người dùng.
\end{dinhnghia}

Hai đặc điểm phân biệt tác tử với một mô hình ngôn ngữ đứng riêng:
\begin{itemize}
  \item \textbf{Tính chủ động hướng mục tiêu.} Người dùng giao \emph{kết quả cần đạt} chứ không giao từng bước; tác tử tự
        quyết định trình tự hành động.
  \item \textbf{Khả năng tác động vào môi trường.} Tác tử không chỉ sinh văn bản mà còn đọc, ghi dữ liệu, gọi API và
        quan sát kết quả hành động của chính nó để điều chỉnh bước tiếp theo.
\end{itemize}

\section{Ba trụ cột kỹ thuật của một tác tử thông minh}
Tổng hợp từ Google Cloud và McKinsey \cite{google_agents,mckinsey_agents}, một tác tử tiêu chuẩn được dựng trên ba trụ
cột quanh lõi mô hình nền tảng (Hình~\ref{fig:ch1-architecture}).

\subsection{Suy luận và lập kế hoạch (Reasoning \& Planning)}
Tác tử phân rã mục tiêu lớn thành chuỗi tác vụ con có thứ tự, đánh giá các phương án và chọn trình tự hợp lý nhất. Mô
hình nền tảng đóng vai bộ não suy luận: hiểu yêu cầu, lập kế hoạch, và khi kết quả trung gian không như mong đợi thì lập
lại kế hoạch.

\subsection{Bộ nhớ và quản lý ngữ cảnh (Memory Management)}
\begin{itemize}
  \item \textbf{Bộ nhớ ngắn hạn} giữ lịch sử hội thoại và trạng thái của phiên làm việc hiện tại, để tác tử hiểu các tham
        chiếu như ``kết quả vừa rồi''.
  \item \textbf{Bộ nhớ dài hạn} lưu tri thức và kinh nghiệm qua nhiều phiên, thường đặt trong cơ sở dữ liệu vector
        (Vector DB) hoặc đồ thị tri thức (Knowledge Graph) và được truy xuất khi cần.
\end{itemize}

\subsection{Sử dụng công cụ ngoại vi (Tool Use \& Action Engine)}
Đây là trụ cột biến ``lời nói'' thành ``hành động''. Tác tử kích hoạt Function Calling để gọi API, truy vấn cơ sở dữ liệu
nội bộ, chạy mã lệnh, tìm kiếm web. Kết quả công cụ trả về trở thành quan sát mới cho vòng suy luận tiếp theo.

\hinhtikz{ch1_architecture}{Sơ đồ kiến trúc nhận thức của một hệ tác tử thông minh (AI Agent Cognitive Architecture):
lõi mô hình nền tảng, ba trụ cột suy luận -- bộ nhớ -- công cụ, môi trường doanh nghiệp và lớp quản trị.}{ch1-architecture}

\section{Vòng tương tác giữa ba trụ cột}
Ba trụ cột không hoạt động rời rạc. Khi nhận mục tiêu, tác tử dùng trụ cột suy luận để lập kế hoạch; mỗi bước của kế
hoạch có thể kích hoạt một công cụ; kết quả công cụ được ghi vào bộ nhớ ngắn hạn và đưa trở lại cho bước suy luận kế
tiếp; tri thức đáng giữ được chuyển vào bộ nhớ dài hạn. Bảng~\ref{tab:ch1-pillars} tóm tắt vai trò của từng trụ cột.

\begin{table}[!htbp]\centering
\caption{Ba trụ cột kỹ thuật của AI Agent và vai trò của chúng.}\label{tab:ch1-pillars}
\small
\begin{tabularx}{\linewidth}{L{3.6cm}XX}
\toprule
\textbf{Trụ cột} & \textbf{Chức năng} & \textbf{Thành phần kỹ thuật điển hình} \\
\midrule
Suy luận \& lập kế hoạch & Phân rã mục tiêu, sắp xếp tác vụ con, lập lại kế hoạch khi gặp lỗi & LLM, kỹ thuật
chuỗi suy nghĩ, mô hình ReAct/ReWOO \\
Bộ nhớ \& ngữ cảnh & Duy trì hội thoại, tích lũy kinh nghiệm qua nhiều phiên & Cửa sổ ngữ cảnh, Vector DB,
Knowledge Graph \\
Công cụ \& hành động & Thu thập dữ liệu thực tế và tác động vào hệ thống & Function Calling, API, truy vấn SQL,
trình thông dịch mã \\
\bottomrule
\end{tabularx}
\end{table}
